\chapter{Ensayos y resultados}
\label{Chapter4}

En este capítulo se presentan las pruebas realizadas al chatbot desarrollado, junto con los resultados obtenidos en distintos escenarios, los procesos de optimización aplicados y un análisis detallado del comportamiento observado. Asimismo, se identifican fallos relevantes y se proponen mejoras para futuras iteraciones del sistema.

Se describen los diferentes escenarios de prueba utilizados para evaluar el rendimiento del chatbot, destacando que la estructura de las pruebas fue clave para detectar áreas susceptibles de mejora y optimización. Además, se analizan los resultados obtenidos en cada escenario con el fin de evaluar la efectividad de las mejoras implementadas, así como identificar posibles limitaciones y proponer ajustes adicionales al sistema.

\section{Evaluación de escenarios del chatbot}
\label{sec:evaluacionEscenarios}

\subsection{Escenario 1: Interacción básica}
Se evaluó la capacidad del chatbot para mantener una conversación coherente en interacciones simples. Para ello, se formularon preguntas directamente relacionadas con el contexto actual de la versión, aunque no siempre completas ni contextualizadas en el hilo conversacional. El objetivo principal de esta prueba fue verificar el denominado \textit{happy path} del sistema; es decir, el funcionamiento correcto de las etapas del \textit{pipeline}, incluyendo la recepción de entrada, la captura de contextos relevantes y el análisis por parte del modelo de lenguaje para generar una respuesta adecuada.

\subsection{Escenario 2: Consultas con múltiples pasos}
En esta evaluación se realizaron consultas más complejas desde el punto de vista del hilo conversacional, donde el usuario planteó preguntas encadenadas que requerían del sistema la retención y gestión del contexto. Se analizó la capacidad del chatbot para mantener coherencia y relevancia a lo largo de múltiples interacciones. Si bien las preguntas fueron formuladas de forma directa, se buscó evaluar la habilidad del sistema para adaptarse a cambios en la temática sin perder el hilo de la conversación ni mezclar las respuestas con temas tratados previamente.

\subsection{Escenario 3: Consultas fuera de contexto}
En este escenario se examinó la capacidad del chatbot para gestionar preguntas ambiguas o fuera de contexto, así como su habilidad para solicitar mayor información o derivar la consulta a un humano. Cabe destacar que esta función se encuentra en etapa de pruebas, ya que el prototipo no cuenta con una conexión real a un sistema CRM para realizar la transferencia de la conversación. El objetivo fue evaluar la habilidad del sistema para reconocer sus propias limitaciones y ofrecer respuestas apropiadas incluso ante la falta de información contextual.


\section{Resultados de las pruebas}
\label{sec:resultadosPruebas}

En este apartado se presentan los resultados obtenidos en cada uno de los escenarios de prueba descritos anteriormente para cada una de las 4 versiones del chatbot. Además, se discuten los hallazgos más relevantes y se proponen recomendaciones para futuras mejoras.

\subsection{Pruebas para la version 0}
\label{sec:resultadosPruebasV0}
\subsubsection{Pruebas de interacción básica}

 

\subsubsection{Pruebas de consultas con múltiples pasos}


\subsubsection{Pruebas de consultas fuera de contexto}


\subsubsection{Conclusiones de la version 0}


\subsection{Pruebas para la version 1}
\label{sec:resultadosPruebasV1}
\subsubsection{Pruebas de interacción básica}

\subsubsection{Pruebas de consultas con múltiples pasos}

\subsubsection{Pruebas de consultas fuera de contexto}

\subsubsection{Conclusiones de la version 1}

\subsection{Pruebas para la version 2}
\label{sec:resultadosPruebasV2}
\subsubsection{Pruebas de interacción básica}

\subsubsection{Pruebas de consultas con múltiples pasos}

\subsubsection{Pruebas de consultas fuera de contexto}

\subsubsection{Conclusiones de la version 2}

\subsection{Pruebas para la version 3}
\label{sec:resultadosPruebasV3}
\subsubsection{Pruebas de interacción básica}

\subsubsection{Pruebas de consultas con múltiples pasos}

\subsubsection{Pruebas de consultas fuera de contexto}

\subsubsection{Conclusiones de la version 3}


\section{Mejoras y optimizaciones propuestas}
\label{sec:mejorasOptimizaciones}

En esta sección se presentan las mejoras y optimizaciones propuestas para el chatbot, basadas en los resultados obtenidos en las pruebas realizadas. Estas recomendaciones están orientadas a abordar los problemas identificados y mejorar la experiencia del usuario, así como la eficiencia del sistema en futuras iteraciones.

\subsection{Mejoras en la gestión del contexto}
\label{sec:mejorasGestionContexto}

\subsection{Optimización del modelo de lenguaje}
\label{sec:optimizacionModeloLenguaje}

\subsection{Integración con sistemas externos}
\label{sec:integracionSistemasExternos}

\subsection{Mejoras en la interfaz de usuario}
\label{sec:mejorasInterfazUsuario}



